%% =====================================================================
\section{Threat Model and Analytical Model}
\label{sec:threat}
%% =====================================================================

\subsection{System Model and Trust Boundaries}
The target is a conventional three-tier web service: a public reverse proxy (Nginx),
a stateless application tier (Flask), an in-memory cache (Redis), and a relational
database (PostgreSQL). Only the proxy is reachable from outside; the application,
cache, and database communicate on a private container network. The product read
endpoint \texttt{GET /products/<id>} follows the cache-aside pattern with a short TTL.
The single trust boundary that matters for this attack is the public interface of
the proxy: everything behind it is assumed to be configured correctly and free of
exploitable defects.

\subsection{Assets and Security Goals}
The protected assets are (a)~the \emph{availability} of the database and (b)~the
\emph{responsiveness} of the service as experienced by legitimate users. We state four
security goals against which the defenses are evaluated:
\begin{description}[leftmargin=1.2em,style=nextline]
\item[G1 -- Database protection.] Under attack, the rate of database work must stay
within the database's capacity and close to its normal level.
\item[G2 -- User availability.] Under attack, legitimate requests must continue to
succeed with latency comparable to normal operation.
\item[G3 -- Timely detection.] An operator must be alerted within seconds of attack
onset, with negligible false alarms during normal operation.
\item[G4 -- Fast recovery.] When defenses are enabled (or the attack stops), the system
must return to a healthy state within seconds, without manual intervention.
\end{description}

\subsection{Adversary Model}
Table~\ref{tab:adversary} lists the adversary's assumed capabilities. The adversary is
an external client who can send any number of well-formed HTTP requests to the public
endpoint and can observe response times. It cannot reach the cache or the database,
cannot read internal metrics, does not exploit software vulnerabilities, and does not
need to forge or poison anything.

\begin{table}[t]
\centering
\caption{Adversary capabilities and assumptions.}
\label{tab:adversary}
\small
\begin{tabularx}{\linewidth}{@{}L{4.0cm}L{1.6cm}Y@{}}
\toprule
Capability & Assumed? & Notes \\
\midrule
Send well-formed HTTP GETs to the public endpoint & Yes & At hundreds of concurrent requests; from one or several source addresses \\
Observe response latency & Yes & Enables a timing side channel (Sec.~\ref{sec:recon}) \\
Know or infer the set of hot keys & Yes & Hot products are typically public (front page, promotions) \\
Know or infer the TTL & Inferable & By probing hit/miss latency over time \\
Reach Redis or PostgreSQL directly & No & Private network \\
Read internal metrics or dashboards & No & Control plane not exposed \\
Exploit code defects, inject, or poison caches & No & Out of scope; attack uses only correct application logic \\
Spoof client identity to the rate limiter & No (by design) & Proxy overwrites \texttt{X-Forwarded-For} (Sec.~\ref{sec:design-rl}) \\
\bottomrule
\end{tabularx}
\end{table}

\subsection{Reconnaissance via a Timing Side Channel}
\label{sec:recon}
The difference between a cache hit and a miss is large and externally visible. In our
testbed the first request for a product is served from the database in
${\approx}137$--$180$~ms, while an immediate repeat is served from Redis in
${\approx}9$--$12$~ms (Section~\ref{sec:sanity}). An adversary can therefore probe a
key periodically and record when responses switch from fast to slow, recovering both
which keys are cached and the TTL. Such probing is low-rate and indistinguishable from
ordinary browsing. This observation matters for defenders: obscuring the TTL is not a
viable mitigation, and TTL jitter only adds noise to an estimate that the attacker can
refine by repeated probing.

\subsection{Attack Chain}
The attack proceeds as follows:
\begin{enumerate}[leftmargin=*]
\item \textbf{Select targets.} Choose a pool of $H$ hot keys whose regeneration is
expensive.
\item \textbf{Synchronize with expiry.} Send a burst of $B$ concurrent requests for one
key at (or after) its expiry. All requests that arrive before the first successful
regeneration miss the cache.
\item \textbf{Fan out to the database.} Each miss issues its own database query; with
a fail-fast pool, misses beyond the pool size fail immediately with HTTP~503, and
because failed regenerations do not populate the cache, subsequent requests in the
burst miss as well.
\item \textbf{Rotate.} Move to the next key in the pool. If the rotation period
exceeds the TTL, every burst lands on an expired entry and the overload becomes
\emph{sustained} rather than a momentary spike.
\end{enumerate}

\subsection{Mapping to Threat Frameworks}
Table~\ref{tab:frameworks} maps the attack onto common threat-modeling vocabularies.
Under STRIDE \cite{shostack2014} it is a pure Denial-of-Service threat. In MITRE
ATT\&CK it falls under Endpoint Denial of Service (T1499), most closely the
\emph{Application Exhaustion Flood} and \emph{Service Exhaustion Flood} sub-techniques
\cite{mitret1499}. In the OWASP API Security Top~10 it is an instance of
\emph{Unrestricted Resource Consumption} (API4:2023) \cite{owaspapi2023}, with the
twist that the consumed resource is downstream of the API rather than the API itself.
From an incident-handling standpoint \cite{nist80061}, the detection subsystem
supports the \emph{detection and analysis} phase and the live defense toggles support
\emph{containment}.

\begin{table}[t]
\centering
\caption{Mapping of the cache-stampede IDDoS to threat-modeling frameworks.}
\label{tab:frameworks}
\small
\begin{tabularx}{\linewidth}{@{}L{3.2cm}L{4.4cm}Y@{}}
\toprule
Framework & Category & Remark \\
\midrule
STRIDE \cite{shostack2014} & Denial of Service & No spoofing, tampering, disclosure, or elevation needed \\
MITRE ATT\&CK \cite{mitret1499} & T1499 Endpoint DoS (Application / Service Exhaustion Flood) & Exhausts the database connection pool via the application \\
OWASP API Top 10 (2023) \cite{owaspapi2023} & API4: Unrestricted Resource Consumption & Resource consumed is downstream (database), not the API process \\
NIST SP 800-61 \cite{nist80061} & Detection \& analysis; containment & Rule/ML alerts; live defense toggles \\
DDoS taxonomy \cite{mirkovic2004} & Application-layer, semantic, non-spoofed & Attack rate can be modest; impact is amplified internally \\
\bottomrule
\end{tabularx}
\end{table}

\subsection{An Analytical Model of the Testbed}
\label{sec:model}
We now derive simple predictions for the testbed's behavior from its parameters
(Table~\ref{tab:params} in Section~\ref{sec:design}). The notation is: $K=100$ products,
TTL $T=4$~s, simulated query time $d_0=0.11$~s, pool size $c=12$, attack burst size
$B=250$, hot-key pool $H=20$, and request rate $\lambda$.

\paragraph{Effective normal request rate.} The normal-traffic generator issues
batches of four concurrent requests and then sleeps $0.1$~s. Its rate is therefore
not the configured nominal value but
$\lambda \approx 4/(0.1 + \ell_{\max})$, where $\ell_{\max}$ is the slowest response in
the batch. With $\ell_{\max}\approx0.12$~s this gives $\lambda\approx18$~req/s,
matching the measured median of 17--18~req/s (Section~\ref{sec:baseline}).

\paragraph{Baseline database load.} Under normal traffic each of the $K$ keys receives
requests at rate $\lambda/K$, approximately Poisson because product ids are drawn
uniformly at random. After a miss refills a key, the key stays fresh for $T$; the next
request after that is a miss. The expected cycle between misses for one key is
therefore $T + K/\lambda$, and the aggregate database rate and miss ratio are
\begin{equation}
q_{\mathrm{base}} = \frac{K}{T + K/\lambda}, \qquad
m_{\mathrm{base}} = \frac{q_{\mathrm{base}}}{\lambda} = \frac{1}{1+\lambda T/K}.
\label{eq:baseline}
\end{equation}
For $\lambda=17$~req/s these give $q_{\mathrm{base}}=10.1$~queries/s and
$m_{\mathrm{base}}=0.60$. A consequence that later proves important is that, with a
short TTL and a uniform catalog, \emph{the normal miss ratio is already high}: most
normal requests are misses because most keys are cold most of the time.

\paragraph{Undefended attack.} A burst of $B$ concurrent requests for an expired key
produces up to $B$ misses. The burst cycle time $\tau_b$ (time for the burst to
complete plus a 50~ms pause) is of the order of the per-request latency under load; at
a measured attack request rate of ${\approx}250$~req/s with $B=250$, $\tau_b\approx1$~s.
Each hot key is therefore revisited every $H\tau_b\approx20$~s, which is much longer
than $T=4$~s, so \emph{every} burst lands on an expired entry. The offered rate of
database attempts is bounded above by the attack request rate, while the nominal pool
capacity, from Eq.~\eqref{eq:capacity}, is
\begin{equation}
\mu_{\max} = c/d_0 = 12/0.11 \approx 109~\text{queries/s}.
\label{eq:mu}
\end{equation}
The model predicts overload whenever the attempt rate exceeds $\mu_{\max}$, and, for a
fail-fast pool, predicts that the excess appears as immediate HTTP~503 failures rather
than as queueing delay. Section~\ref{sec:mechanics} shows that the real system
saturates much earlier than Eq.~\eqref{eq:mu} suggests.

\paragraph{Defended attack.} With single-flight locking, at most one regeneration per
key can be in progress, and while a fresh entry exists none is needed. The database
rate caused by hot keys is therefore bounded by one query per key per TTL,
\begin{equation}
q_{\mathrm{hot}} \le \frac{H}{T} = \frac{20}{4} = 5~\text{queries/s},
\label{eq:defended}
\end{equation}
independently of $B$, and the total across the whole catalog is bounded by
$K/T=25$~queries/s. The edge limiter additionally caps the request rate reaching the
application from one source address at its configured rate $r=50$~req/s plus a burst
allowance of 20. The model thus predicts a defended database load of a few queries per
second, a backend request rate pinned at ${\approx}50$~req/s, and a \emph{coalescing
ratio}---cache misses per database query---well above one.

\paragraph{Amplification.} We define the stampede amplification of an expiry event as
the number of database queries it causes, $A=Q_{\mathrm{exp}}$. Undefended, $A$ is
bounded only by the burst size ($A\le B=250$); defended, $A=1$. The attacker's leverage
is the ratio of database work caused to requests sent. Section~\ref{sec:results}
reports both quantities as measured.
